\subsection{Revisão Documental por Docentes}
\label{sec:revisao-docente-operacional}

Propõe-se um painel de cinco docentes distintos, selecionados intencionalmente: três licenciados em matemática com pelo menos dois anos de docência na educação básica e dois com atuação em ensino superior ou pesquisa em educação matemática, também com pelo menos dois anos de experiência. Ao menos um deverá ter experiência em avaliação educacional ou tecnologias educacionais. Os perfis poderão se sobrepor, mas cada pessoa contará uma vez. O autor e o orientador não comporão esse painel; organizarão a documentação e registrarão as decisões. Trata-se de um plano de execução, sem participantes convidados ou avaliação realizada.

Cada docente receberá a mesma versão da especificação, a matriz de origem dos requisitos, os cenários e o formulário do Apêndice~\ref{apendice:formulario-oe4}. A primeira rodada será individual, com cerca de 30 minutos de leitura e 30 de preenchimento. Para cada R01--R26, serão julgadas quatro dimensões: sustentação na fonte ou decisão autoral, adequação curricular, clareza da saída para o professor e viabilidade documental das condições necessárias. A escala proposta é: 1=incoerente; 2=requer alteração substancial; 3=coerente com ajuste pontual; 4=coerente. A opção NE significa evidência insuficiente e exige justificar o que falta. O campo de aceitação funcional relacionará F01--F07 ao critério correspondente do Quadro~\ref{qua:aceitacao-funcional}.

Um requisito será considerado ``coerente'' se pelo menos quatro dos cinco docentes lhe atribuírem 3 ou 4 em cada dimensão e não houver problema crítico pendente. Serão críticos, para este protocolo, uma saída que trate previsão como proficiência sem mensuração validada, uma prescrição sem revisão docente ou exposição indevida de dados. Nos demais casos, o resultado será ``revisar'', com indicação do trecho e da alteração necessária. Essa regra de 80\% é uma escolha de projeto para o painel; não constitui validação estatística do instrumento nem estimativa para uma população docente.

Após a primeira rodada, o autor reunirá as divergências e proporá uma versão revista, preservando as respostas originais. Uma reunião de até 60 minutos permitirá discutir os motivos, sem substituir os julgamentos individuais por concordância forçada. Haverá uma segunda rodada individual, com a mesma regra. Persistindo divergência que impeça atingir a regra de consenso ou problema crítico, o requisito permanecerá ``revisar''. Ajustes pontuais associados à resposta 3 deverão ser incorporados à versão final e conferidos pelos docentes que os solicitaram; até essa confirmação, a aceitação ficará condicionada ao ajuste. Alteração substancial exigirá manter ``revisar''. A especificação será classificada como coerente documentalmente apenas se todos os requisitos atenderem à regra e os ajustes estiverem confirmados; caso contrário, será apresentada com a lista dos que exigem revisão. Com menos de cinco avaliações completas, não se aplicará a regra de consenso: o painel será registrado como incompleto. A conferência técnica de dados, entradas e saídas será registrada separadamente, sem ser usada como voto docente.

Antes de recrutar participantes ou coletar respostas, o pesquisador responsável e o orientador deverão encaminhar o projeto e seus instrumentos ao comitê de ética competente, pela Plataforma Brasil, e aguardar a decisão aplicável. Serão previstos consentimento livre e esclarecido, participação voluntária, possibilidade de desistência, riscos de constrangimento profissional e sua redução, acesso restrito às respostas e divulgação sem identificação dos docentes \cite[art.~3, 17, 23 e 28]{CNS5102016}. O IFC disponibiliza orientações e modelos para a submissão e o consentimento.\footnote{CEPSH/IFC, Formulário e Documentos: \url{https://cepsh.ifc.edu.br/documentos/}. Acesso em: 29 out. 2026.} A avaliação utilizará exemplos fictícios, sem dados reais de estudantes, e não avaliará o desempenho profissional dos participantes. A aplicação desse plano permanece posterior à entrega documental e à tramitação ética.
